%% ------------------------------------------------------------------
%% identity.tex — Shared visual identity for CV and cover letters
%%
%% \input this file from any LaTeX document that should share the
%% Gonçalves visual identity (colors, icons, tipography helpers,
%% graphics search paths).
%%
%% Idempotent by design: safe to \input from documents that already
%% loaded xcolor / fontawesome / graphicx (via altacv or otherwise).
%% Uses \providecommand for macros so re-\input never errors.
%%
%% Callers currently include:
%%   - applications/_shared/latex/cover-letter-template.tex
%%   - applications/YYYY-MM-<slug>/cover-letter.tex (copied from template)
%%   - Optionally the CV itself, if the shared identity refactor is applied.
%% ------------------------------------------------------------------

%% ------- Required packages (loaded only if not already loaded) -------
\makeatletter
\@ifpackageloaded{xcolor}{}{\usepackage{xcolor}}
\@ifpackageloaded{fontawesome}{}{\usepackage{fontawesome}}
\@ifpackageloaded{graphicx}{}{\usepackage{graphicx}}
\@ifpackageloaded{tikz}{}{\usepackage{tikz}}
\makeatother

%% ------- Color palette (same values used by curriculum-en.tex) -------
\definecolor{StPatricksBlue}{HTML}{1F3D77}
\definecolor{SlateGrey}{HTML}{2E2E2E}
\definecolor{LightGrey}{HTML}{333333}

%% ------- Graphics search paths -------
%% Lets \includegraphics{profile} resolve to english/profile.jpg from any
%% depth the caller might be at. LaTeX tries these in order and uses the
%% first match.
\graphicspath{%
    {./}%                          same folder as the caller
    {./english/}%                  from project root
    {../english/}%                 from applications/ (unlikely, but cheap)
    {../../english/}%              from applications/YYYY-MM-<slug>/
    {../../../english/}%           from applications/_shared/latex/
}

%% ------- Style helpers for cover letters -------
%% Not used by the CV (altacv defines its own header macros), but harmless
%% to have loaded. \providecommand only defines if not already defined.

\providecommand{\senderName}[1]{%
    {\fontsize{22pt}{28pt}\selectfont\bfseries\color{StPatricksBlue}#1}%
}

\providecommand{\senderTagline}[1]{%
    {\fontsize{11pt}{14pt}\selectfont\itshape\color{SlateGrey}#1}%
}

\providecommand{\contactBullet}{%
    \hspace{0.5em}\textcolor{StPatricksBlue}{$\bullet$}\hspace{0.5em}%
}

\providecommand{\dividerLine}{%
    {\color{StPatricksBlue}\rule{\linewidth}{0.6pt}}%
}

\providecommand{\salutation}[1]{%
    {\bfseries\color{SlateGrey}#1}%
}

\providecommand{\closing}[1]{%
    {\bfseries\color{SlateGrey}#1}%
}

%% ------- Date and subject-line helpers for cover letters -------

%% \letterDate{City}
%% Renders "City, DD Month YYYY" right-aligned, in italic SlateGrey. Uses
%% \today under the hood so the date is always current at compile time.
%% Example: \letterDate{Lausanne}  ->  "Lausanne, 27 July 2026"
\providecommand{\letterMonth}{%
    \ifcase\month\or January\or February\or March\or April\or May\or June%
    \or July\or August\or September\or October\or November\or December\fi%
}
\providecommand{\letterDate}[1]{%
    \hfill{\itshape\color{SlateGrey}#1,~\the\day~\letterMonth~\the\year}%
}

%% \subjectLine{content}
%% Renders a bold SlateGrey subject line. Common convention in formal
%% European/Swiss/French letters (e.g. "Subject: ..." or "Object: ...").
%% The caller decides whether to include a prefix.
%% Example: \subjectLine{Subject: Application for Full Stack Engineer -- Ref. 2357}
\providecommand{\subjectLine}[1]{%
    {\bfseries\color{SlateGrey}#1}%
}

%% ------- Circular photo (matches altacv's \photo behavior in the CV) -------
%% \circularPhoto{diameter}{filename}
%% Renders a circular-cropped photo, resolving the filename via \graphicspath
%% (so \circularPhoto{2.5cm}{profile} finds english/profile.jpg).
%% The tikzpicture baseline is set to the top of the circle so that when used
%% inside a minipage with [t] alignment, the picture aligns with the top of
%% the row rather than the middle of the circle (default tikz behavior).
%% Example: \circularPhoto{2.5cm}{profile}
\providecommand{\circularPhoto}[2]{%
    \begin{tikzpicture}[baseline=(currentphoto.north)]
        \node[circle,minimum size=#1,inner sep=0pt,%
            path picture={%
                \node[inner sep=0pt] at (path picture bounding box.center) {%
                    \includegraphics[width=#1]{#2}%
                };%
            }%
        ] (currentphoto) {};
    \end{tikzpicture}%
}
